\addcontentsline{toc}{subsubsection}{Solution 20. The Cross-Dilated Shift}

\begin{dlsummary}[title={\bfseries Solution 20. The Cross-Dilated Shift},breakable]

\begin{proof}
The identity function satisfies the inequality with equality. We prove it
is the only solution.

\medskip
\noindent
\textbf{Step 1. Establish a lower bound.}

Since $f(xf(y))>0$, the given inequality implies $f(x+y)>x$ for all
$x,y>0$. Fix $z>0$. For each $\varepsilon\in(0,z)$, take
$x=z-\varepsilon$ and $y=\varepsilon$. Then $f(z)>z-\varepsilon$;
hence
\[
f(z)\geq z\qquad(z>0). \tag{1}
\]

\medskip
\noindent
\textbf{Step 2. Bound the excess over the identity.}

Suppose $f(a)>a+1$ for some $a>0$, and set
\[
x=\frac{a}{f(a)-1}.
\]
Then $0<x<1$ and $xf(a)=x+a$. Applying the given inequality at
$(x,a)$ gives
\[
x^2+f(x+a)\leq x f(x+a),
\]
which is impossible because $f(x+a)>0$ and $0<x<1$. Therefore
$f(a)\leq a+1$ for all $a>0$.

Define $g(t)=f(t)-t$. By (1),
\[
0\leq g(t)\leq1\qquad(t>0). \tag{2}
\]
Substituting $f(t)=t+g(t)$ into the original inequality and cancelling
the common terms yields
\[
xg(y)+g\bigl(xf(y)\bigr)\leq xg(x+y)
\qquad(x,y>0). \tag{3}
\]

\medskip
\noindent
\textbf{Step 3. Force the excess to vanish.}

Since $g$ is nonnegative, (3) implies $g(y)\leq g(x+y)$ for all
$x,y>0$, so $g$ is nondecreasing. Taking $x=1$ in (3), and using
$f(y)\geq y$ and the monotonicity of $g$, gives
\[
2g(y)\leq g(y)+g\bigl(f(y)\bigr)\leq g(y+1).
\]
Inductively, $g(y+n)\geq2^n g(y)$ for every positive integer $n$.
By (2), this forces $g(y)=0$ for every $y>0$. Thus $f(y)=y$ on
$(0,\infty)$, as claimed.
\end{proof}

\end{dlsummary}